\chapter{Methodology: Job-Centric MARL Framework}
\label{chap:methodology}

This chapter presents the proposed job-centric Multi-Agent Reinforcement Learning (MARL) framework developed for solving the Flexible Job-Shop Scheduling Problem (FJSP) under stochastic job arrivals.
The approach combines a Discrete Event Simulation (DES) environment with a MARL algorithm (MASA-QMIX) to enable adaptive, cooperative, and decentralized scheduling decisions in dynamic production settings.

\section{Problem Definition}
\label{sec:problemdefinition}
This research addresses a dual-resource Flexible Job-Shop Scheduling Problem (FJSP) under stochastic job arrivals, where production efficiency depends on coordinated real-time allocation of both machines and operators in a dynamic environment.

The job shop system consists of \( M \) machines, \( K \) operators, and a continuously evolving set of heterogeneous jobs \( J \). Each job \( j \in J \) comprises a predefined sequence of operations \( O_j = \{o_{j,1}, o_{j,2}, \ldots, o_{j,n_j}\} \), where each operation must be executed on a compatible machine-operator pair before the next operation can begin.

\textbf{ Key Problem Characteristics:}
\begin{itemize}
\item \textbf{Dual-Resource Allocation:}
Each operation requires simultaneous assignment of both a qualified machine and a qualified operator. The scheduling decision must select a feasible (machine, operator) pair from the set of compatible resource combinations, where compatibility is determined by machine capability and operator skill qualifications.

\item \textbf{Stochastic Job Arrivals and Dynamic Agent Insertion:}
New jobs enter the system during execution following a stochastic arrival process. Unlike static scheduling scenarios, the scheduler operates without knowledge of future arrivals and must adapt policies in real-time as the workload evolves. Each arriving job instantiates a new autonomous agent that immediately participates in the learned cooperative policy without restarting the simulation or retraining the model.

\item \textbf{Heterogeneous Jobs:}
Jobs differ in operation sequences, processing time requirements, and machine-operator qualification constraints. Processing times for the same operation may vary depending on which machine-operator pair executes it, reflecting real-world skill and capability differences.

\item \textbf{Operational Constraints:}
The scheduling environment operates under the following constraints:

\begin{itemize}
    \item Each machine processes at most one operation at a time (no parallel execution on a single machine) — \textit{\textbf{(machine disjunction constraint)}},
    \item Each operator can be assigned to at most one machine-operation pair at any given time — \textit{\textbf{(operator disjunction constraint)}},
    \item Operations within a job must execute in sequential order (operation \( o_{j,i+1} \) cannot start until \( o_{j,i} \) completes) — \textit{\textbf{(precedence constraint)}},
    \item Operations cannot be preempted once started — \textit{\textbf{(non-preemption constraint)}},
    \item Operations can only execute on (machine, operator) pairs where both the machine is capable of the operation type and the operator is qualified for that machine — \textit{\textbf{(capability constraints)}},
    \item Jobs unable to find an available feasible (machine, operator) pair wait until the next decision point and attempt allocation again.
\end{itemize}
\end{itemize}
\textbf{Optimization Objectives:}
The scheduling policy aims to:
\begin{itemize}
    \item Minimize job waiting times (average flow time),
    \item Minimize total makespan,
    \item Maximize resource utilization and load balance,
    \item Maximize system throughput (jobs completed per time unit).
\end{itemize}

\textbf{Multi-Agent Reinforcement Learning Approach:}
This work formulates scheduling as a Multi-Agent Reinforcement Learning (MARL) problem where each job acts as an autonomous agent responsible for selecting its own resource allocations. Agents learn cooperative policies through shared rewards without requiring centralized coordination during execution. This job-centric design naturally handles dynamic agent populations and enables decentralized decision-making under partial observability-essential properties for real-time scheduling in stochastic production environments.

\section{Formal Problem Formulation (Dec-POMDP / Stochastic Game)}
\label{sec:problemformulation}

The dual-resource FJSP with stochastic job arrivals is formulated as a Decentralized Partially Observable Markov Decision Process (Dec-POMDP), which provides a rigorous mathematical framework for multi-agent decision-making under partial observability and uncertainty. This formulation naturally captures the decentralized nature of job-centric scheduling, where each job agent makes resource allocation decisions based on local observations while coordinating implicitly through shared system dynamics and rewards.

\subsection{Mathematical Framework}
A Dec-POMDP is defined by the tuple \( \langle \mathcal{N}, \mathcal{S}, \{\mathcal{A}^i\}, \mathcal{T}, \{\mathcal{O}^i\}, \mathcal{Z}, \mathcal{R}, \gamma \rangle \), where:

\begin{itemize}
    \item \( \mathcal{N} = \{1, 2, \ldots, n(t)\} \) is the set of agents (jobs) at time \( t \), where \( n(t) \) varies dynamically as jobs arrive and complete,
    
    \item \( \mathcal{S} \in \mathbb{R}^{d_s} \) is the global state space encoding system-wide metrics,
    
    \item \( \mathcal{A}^i \subseteq \{1, \ldots, M\} \) is the action space for agent \( i \), representing machine selection,
    
    \item \( \mathcal{T}: \mathcal{S} \times \mathcal{A}^1 \times \cdots \times \mathcal{A}^n \rightarrow \Delta(\mathcal{S}) \) is the state transition function,
    
    \item \( \mathcal{O}^i \in \mathbb{R}^{d_o} \) is the observation space for agent \( i \),
    
    \item \( \mathcal{Z}: \mathcal{S} \times \mathcal{A}^1 \times \cdots \times \mathcal{A}^n \rightarrow \Delta(\mathcal{O}^1 \times \cdots \times \mathcal{O}^n) \) is the observation function,
    
    \item \( \mathcal{R}: \mathcal{S} \times \mathcal{A}^1 \times \cdots \times \mathcal{A}^n \rightarrow \mathbb{R} \) is the shared reward function,
    
    \item \( \gamma \in [0,1) \) is the discount factor for future rewards.
\end{itemize}

The objective is to learn a joint policy \( \boldsymbol{\pi} = (\pi^1, \ldots, \pi^n) \) that maximizes the expected cumulative discounted reward:
\[
J(\boldsymbol{\pi}) = \mathbb{E}_{\tau \sim \boldsymbol{\pi}} \left[ \sum_{t=0}^{T} \gamma^t r_t \right]
\]
where \( \tau = (s_0, \mathbf{a}_0, r_0, s_1, \mathbf{a}_1, r_1, \ldots) \) denotes a trajectory under policy \( \boldsymbol{\pi} \).

\subsection{State Space $\mathcal{S}$: Global System Representation}
The global state \( s_t \in \mathbb{R}^{10} \) provides a complete snapshot of the system at decision epoch \( t \):
\begin{equation}
s_t = \begin{bmatrix}
n_{\text{jobs}}^{\text{arrived}} & n_{\text{jobs}}^{\text{proc}} & n_{\text{jobs}}^{\text{wait}} & n_{\text{ops}}^{\text{arrived}} & n_{\text{ops}}^{\text{proc}} \\
n_{\text{ops}}^{\text{wait}} & u_M & u_O & W_{\text{total}} & \tau
\end{bmatrix}^\top
\end{equation}
where:
\begin{itemize}
    \item \( n_{\text{jobs}}^{\text{arrived}}, n_{\text{jobs}}^{\text{proc}}, n_{\text{jobs}}^{\text{wait}} \): raw counts of arrived, processing, and waiting jobs,
    \item \( n_{\text{ops}}^{\text{arrived}}, n_{\text{ops}}^{\text{proc}}, n_{\text{ops}}^{\text{wait}} \): raw counts of arrived, processing, and waiting operations,
    \item \( u_M = \frac{\sum_{m=1}^{M} \mathbb{1}[\text{machine } m \text{ busy}]}{M} \in [0,1] \): normalized machine utilization,
    \item \( u_O = \frac{\sum_{k=1}^{K} \mathbb{1}[\text{operator } k \text{ busy}]}{K} \in [0,1] \): normalized operator utilization,
    \item \( W_{\text{total}} = \sum_{j \in J_{\text{active}}} w_j \): cumulative waiting time (raw, unbounded),
    \item \( \tau = \frac{t_{\text{sim}}}{T_{\text{episode}}} \in [0,1] \): normalized episode time fraction.
\end{itemize}

\subsection{Observation Space $\mathcal{O}^i$: Local Agent View}
Each agent \( i \) receives a partial observation \( o_t^i \in \mathbb{R}^8 \) encoding only its own job-specific information:
\begin{equation}
o_t^i = \begin{bmatrix}
\theta_{\text{current}} & n_{\text{ops}}^{\text{total}} & n_{\text{ops}}^{\text{remain}} & w_i & |M_\theta| & |M_{\theta}^{\text{free}}| & n_{\text{jobs}}^{\text{active}} & \delta_{\text{finished}}
\end{bmatrix}^\top
\end{equation}
where:
\begin{itemize}
    \item \( \theta_{\text{current}} \in \{0, 1, \ldots, 9\} \): operation type index of current operation,
    \item \( n_{\text{ops}}^{\text{total}}, n_{\text{ops}}^{\text{remain}} \): total and remaining operations for this job,
    \item \( w_i \geq 0 \): cumulative waiting time for job \( i \) (raw, unbounded),
    \item \( |M_\theta| \): number of machines theoretically capable of operation type \( \theta \) (theoretical capacity),
    \item \( |M_{\theta}^{\text{free}}| \): number of capable machines currently available after feasibility masking,
    \item \( n_{\text{jobs}}^{\text{active}} \): system-wide active job count (shared coordination signal),
    \item \( \delta_{\text{finished}} \in \{0,1\} \): job completion indicator (1 when all operations complete, 0 otherwise).
\end{itemize}

\textbf{Partial Observability:} Each agent observes only its own job state (\( \theta_{\text{current}}, n_{\text{ops}}^{\text{remain}}, w_i \)) and limited system-level information (\( |M_{\theta}^{\text{free}}|, n_{\text{jobs}}^{\text{active}} \)). Agents cannot observe: (1) detailed schedules of other jobs, (2) exact resource allocation timelines, (3) future job arrivals, or (4) the full global state \( s_t \). This enforces realistic information constraints present in distributed manufacturing systems.

\subsection{Action Space $\mathcal{A}^i$ and Action Masking}
Each agent \( i \)'s action space is a discrete set of machine indices:
\[
\mathcal{A}^i = \{1, 2, \ldots, M\}
\]
where \( M \) is the total number of machines in the system. Although agents formally select machines, \textit{the action masking mechanism ensures that each selectable action implicitly represents a feasible (machine, operator) pair}: agents can only observe and select machines for which both the machine is available \textit{and} a qualified operator exists and is currently free. This implicit pair feasibility is enforced through the three-condition mask validation before agents make decisions.

\textbf{Action Masking:} At each decision epoch, a binary feasibility mask \( \mathcal{M}_t^i \in \{0,1\}^M \) is constructed for agent \( i \):
\begin{equation}
\mathcal{M}_t^i[m] = \mathbb{1}\left[ C_{\theta,m} \land \neg B_m \land \exists k \in \mathcal{K}_m^{\text{free}} : Q_{k,\theta,m} \right]
\end{equation}
where:
\begin{itemize}
    \item \( C_{\theta,m} \): machine \( m \) is capable of operation type \( \theta \) (\textit{capability constraint}),
    \item \( B_m \): machine \( m \) is currently busy processing another operation,
    \item \( \mathcal{K}_m^{\text{free}} \): set of operators currently not assigned to any machine,
    \item \( Q_{k,\theta,m} \): operator \( k \) is qualified for machine \( m \).
\end{itemize}

This three-condition conjunction ensures \textit{dual-resource feasibility}: a machine \( m \) is only selectable (\( \mathcal{M}_t^i[m] = 1 \)) if it is capable, not busy, \textit{and} has at least one qualified operator available. Thus, the masked action space \( \mathcal{A}_t^i = \{ a \in \mathcal{A}^i : \mathcal{M}_t^i[a] = 1 \} \) implicitly represents the set of feasible \textit{(machine, operator)} pairs, where each machine index corresponds to a valid resource combination. If \( \mathcal{A}_t^i = \emptyset \) (either because no capable machines are currently available or because no qualified operator is free for the feasible machines), the agent defers its decision and waits for the next decision epoch.

\textbf{Operator Assignment Coupling:} When agent \( i \) selects machine \( m \) from the masked action space \( \mathcal{A}_t^i \), the environment assigns a qualified operator \( k^* \) from the set of available operators that satisfied the mask condition:
\[
k^* = \arg\min_{k \in \mathcal{K}_m^{\text{free}} : Q_{k,\theta,m}} \, \xi_k, \quad \xi_k \sim \text{Uniform}(0,1)
\]
where the random seed is fixed per episode for reproducibility. Since the action mask already verified \( \exists k \in \mathcal{K}_m^{\text{free}} : Q_{k,\theta,m} \), this assignment is guaranteed to succeed. This design eliminates the combinatorial explosion of explicit \textit{(machine, operator)} action spaces: instead of \( \mathcal{A}^i = \prod_{m=1}^M \prod_{k=1}^K \{(m,k)\} \), the action space remains \( \mathcal{A}^i = \{1,\ldots,M\} \), with masking implicitly encoding pair feasibility. Critically, operators function as an \textit{additional resource constraint} rather than a decision variable.

\subsection{Transition Dynamics $\mathcal{T}$ and Decision Epochs}
State transitions are governed by a hybrid SimPy discrete-event simulation:
\begin{equation}
s_{t+1} \sim \mathcal{T}(s_t, \mathbf{a}_t), \quad \mathbf{a}_t = (a_t^1, \ldots, a_t^{n(t)})
\end{equation}

\textbf{Decision Epoch Mechanism:} Decision epochs occur when jobs reach decision points in their lifecycle:
\begin{enumerate}
    \item Each active job runs a SimPy coroutine (\texttt{\_job\_process}) that yields control when it requires a resource allocation decision.
    \item The job pushes a decision request to a queue (\texttt{pending\_decisions}) and suspends via \texttt{yield resume\_evt}.
    \item The training loop polls the simulation in discrete time steps (\( \Delta t = 1.0 \)) using \texttt{wait\_for\_decisions()} until \texttt{pending\_decisions} is non-empty.
    \item All pending jobs form a decision batch \( \{j_1, \ldots, j_{|\mathcal{B}_t|}\} \), and agents simultaneously select machines via the policy \( \pi_\theta \).
    \item Selected machines are allocated with environment-assigned operators, operations begin execution, and the simulation resumes until the next decision epoch.
\end{enumerate}

\textbf{Processing Time Heterogeneity and Job Completion Variability:}
Operation processing times \( d_{\theta,m} \) exhibit \textit{processing time variability} across different machines due to varying machine capabilities. For a given operation type \( \theta \), the processing time is \textit{deterministically} determined by the selected machine \( m \), reflecting differences in machine efficiency and technological specifications. These durations are derived from empirically-calibrated processing time means \texttt{processing\_time\_means}[\(\theta\)][\(m\)], which capture the expected execution time for each operation-machine combination.

Job completion times are highly \textit{variable} and depend on multiple dynamic factors:
\begin{itemize}
    \item \textbf{Scheduling Decisions:} The sequence and timing of machine-operator allocations directly affect when operations begin and complete.
    \item \textbf{Waiting Times:} Resource contention causes jobs to wait when all feasible machine-operator pairs are busy, accumulating delay that propagates through the job's operation sequence. Since agents cannot directly control operator assignment, operator availability introduces stochastic wait times: a machine may be free but unavailable if no qualified operators are idle.
    \item \textbf{System Load Dynamics:} Stochastic job arrivals following a compound Poisson process \( N(t) \sim \text{Poisson}(\lambda t) \), where \( \lambda \) is the arrival rate parameter, create fluctuating workload conditions. Periods of high arrival density lead to congestion where queue lengths grow and jobs experience extended wait times, while low-arrival periods allow the system to clear backlogs.
    \item \textbf{Dual-Resource Constraint:} Operator availability at each decision epoch depends on current operator utilization. While agents learn to select machines, the environment controls operator assignment, creating an implicit coupling where the same machine selection may yield different wait times depending on operator availability patterns. Agents learn temporal coordination through global operator utilization signals (\( u_O \)) and wait time feedback, adapting machine selection timing to avoid contested resources.
\end{itemize}

Consequently, two identical jobs arriving at different times or scheduled differently can exhibit vastly different completion times, even though individual operation durations \( d_{\theta,m} \) remain constant for each (operation, machine) pair.

Between decision epochs, the simulation advances as operations execute for their prescribed durations \( d_{\theta,m} \), and the global state evolves continuously. New decision epochs are triggered when:
\begin{enumerate}
    \item An operation completes, releasing machine-operator resources and allowing the job to request allocation for its next operation.
    \item A new job arrives stochastically, instantiating a new agent \( i_{\text{new}} \) that immediately enters the decision queue.
\end{enumerate}

\subsection{Reward Function $\mathcal{R}$: Multi-Objective Optimization}
All agents share a global reward signal designed to balance competing scheduling objectives:
\begin{equation}
r_t = w_1 C_{\text{norm}} - w_2 W_{\text{norm}} - w_3 P_{\text{norm}} + w_4 \Delta_{\text{throughput}} + w_5 B_{\text{load}}
\end{equation}
where:
\begin{itemize}
    \item \( C_{\text{norm}} = \frac{|\{j \in J : j \text{ completed}\}|}{|J|} \in [0,1] \): normalized completion rate,
    \item \( W_{\text{norm}} = \frac{\bar{w}}{W_{\max}} \in [0,1] \): normalized average waiting time, \( \bar{w} = \frac{1}{|J_{\text{active}}|} \sum_{j \in J_{\text{active}}} w_j \),
    \item \( P_{\text{norm}} = \frac{|J_{\text{active}}|}{|J|} \in [0,1] \): normalized work-in-progress ratio,
    \item \( \Delta_{\text{throughput}} = \frac{|\{j : j \text{ completed in last step}\}|}{|J|} \): instantaneous throughput change,
    \item \( B_{\text{load}} = \lambda_M \cdot H(M) + \lambda_O \cdot H(O) \in [0,1] \): entropy-based load balance score,
\end{itemize}
with entropy \( H(M) = -\sum_{m=1}^M p_m \log p_m \), where \( p_m = \frac{c_m}{\sum_{m'} c_{m'}} \) and \( c_m \) is the count of recent decisions selecting machine \( m \) over a sliding window of recent decisions.

\textbf{Reward Weight Configuration:}
The weights \( w_1, w_2, w_3, w_4, w_5 \) and parameters \( \lambda_M, \lambda_O, W_{\max} \) are hyperparameters that control the relative importance of each objective component. Higher values of \( w_4 \) provide stronger immediate feedback for throughput improvements, while higher \( w_2 \) increases the penalty for excessive waiting times. The load balance weight \( w_5 \) encourages diverse resource usage to prevent bottlenecks. These weights are tuned during training to balance competing objectives and directly address the optimization goals defined in Section~\ref{sec:problemdefinition}.

\subsection{Policy Representation and Learning Objective}
Each agent \( i \) follows a parameterized policy \( \pi_\theta^i(a^i \mid h_t^i) \), where \( h_t^i \) is the agent's recurrent hidden state encoding its observation history. Under the MASA-QMIX framework (detailed in Section~\ref{sec:masaqmix}), agents share RNN parameters but maintain separate hidden states, enabling:
\begin{itemize}
    \item \textit{Parameter Sharing:} Single RNN \( f_\theta: \mathcal{O}^i \times \mathcal{H} \rightarrow \mathbb{R}^{|\mathcal{A}^i|} \times \mathcal{H} \) processes all agents,
    \item \textit{Dynamic Population Handling:} New agents use the same \( f_\theta \) with zero-initialized \( h_0^i \),
    \item \textit{Action Selection:} \( \epsilon \)-greedy exploration:
    \[
    a_t^i = 
    \begin{cases}
    \text{uniform}(\mathcal{A}_t^i) & \text{with probability } \epsilon(t) \\
    \arg\max_{a \in \mathcal{A}_t^i} Q_\theta(h_t^i, a) & \text{otherwise}
    \end{cases}
    \]
    where \( \epsilon(t) \) decays linearly from \( \epsilon_{\text{start}} \) to \( \epsilon_{\text{end}} \) over the training horizon to balance exploration and exploitation.
\end{itemize}

The learning objective is to maximize the expected return by minimizing the temporal difference error:
\[
\mathcal{L}(\theta) = \mathbb{E}_{\tau \sim \mathcal{D}} \left[ \left( Q_{\text{tot}}(\mathbf{s}_t, \mathbf{a}_t; \theta) - y_t \right)^2 \right]
\]
where \( y_t = r_t + \gamma \max_{\mathbf{a}'} Q_{\text{tot}}(\mathbf{s}_{t+1}, \mathbf{a}'; \theta^-) \) is the target Q-value computed using a separate target network with frozen parameters \( \theta^- \), updated periodically to stabilize learning. The replay buffer \( \mathcal{D} \) stores episode trajectories for experience replay, with mini-batches sampled during training to break temporal correlations and improve sample efficiency.

\section{System Architecture}
\label{sec:architecture}

\chapter{Methodology: Job-Centric MARL Framework}
\label{chap:methodology}

This chapter presents the proposed job-centric Multi-Agent Reinforcement Learning (MARL) framework developed for solving the Flexible Job-Shop Scheduling Problem (FJSP) under stochastic job arrivals.
The approach combines a Discrete Event Simulation (DES) environment with a MARL algorithm (MASA-QMIX) to enable adaptive, cooperative, and decentralized scheduling decisions in dynamic production settings.

\section{Problem Definition}
\label{sec:problemdefinition}
This research addresses a dual-resource Flexible Job-Shop Scheduling Problem (FJSP) under stochastic job arrivals, where production efficiency depends on coordinated real-time allocation of both machines and operators in a dynamic environment.

The job shop system consists of \( M \) machines, \( K \) operators, and a continuously evolving set of heterogeneous jobs \( J \). Each job \( j \in J \) comprises a predefined sequence of operations \( O_j = \{o_{j,1}, o_{j,2}, \ldots, o_{j,n_j}\} \), where each operation must be executed on a compatible machine-operator pair before the next operation can begin.

\textbf{ Key Problem Characteristics:}
\begin{itemize}
\item \textbf{Dual-Resource Allocation:}
Each operation requires simultaneous assignment of both a qualified machine and a qualified operator. The scheduling decision must select a feasible (machine, operator) pair from the set of compatible resource combinations, where compatibility is determined by machine capability and operator skill qualifications.

\item \textbf{Stochastic Job Arrivals and Dynamic Agent Insertion:}
New jobs enter the system during execution following a stochastic arrival process. Unlike static scheduling scenarios, the scheduler operates without knowledge of future arrivals and must adapt policies in real-time as the workload evolves. Each arriving job instantiates a new autonomous agent that immediately participates in the learned cooperative policy without restarting the simulation or retraining the model.

\item \textbf{Heterogeneous Jobs:}
Jobs differ in operation sequences, processing time requirements, and machine-operator qualification constraints. Processing times for the same operation may vary depending on which machine-operator pair executes it, reflecting real-world skill and capability differences.

\item \textbf{Operational Constraints:}
The scheduling environment operates under the following constraints:

\begin{itemize}
    \item Each machine processes at most one operation at a time (no parallel execution on a single machine) — \textit{\textbf{(machine disjunction constraint)}},
    \item Each operator can be assigned to at most one machine-operation pair at any given time — \textit{\textbf{(operator disjunction constraint)}},
    \item Operations within a job must execute in sequential order (operation \( o_{j,i+1} \) cannot start until \( o_{j,i} \) completes) — \textit{\textbf{(precedence constraint)}},
    \item Operations cannot be preempted once started — \textit{\textbf{(non-preemption constraint)}},
    \item Operations can only execute on (machine, operator) pairs where both the machine is capable of the operation type and the operator is qualified for that machine — \textit{\textbf{(capability constraints)}},
    \item Jobs unable to find an available feasible (machine, operator) pair wait until the next decision point and attempt allocation again.
\end{itemize}
\end{itemize}
\textbf{Optimization Objectives:}
The scheduling policy aims to:
\begin{itemize}
    \item Minimize job waiting times (average flow time),
    \item Minimize total makespan,
    \item Maximize resource utilization and load balance,
    \item Maximize system throughput (jobs completed per time unit).
\end{itemize}

\textbf{Multi-Agent Reinforcement Learning Approach:}
This work formulates scheduling as a Multi-Agent Reinforcement Learning (MARL) problem where each job acts as an autonomous agent responsible for selecting its own resource allocations. Agents learn cooperative policies through shared rewards without requiring centralized coordination during execution. This job-centric design naturally handles dynamic agent populations and enables decentralized decision-making under partial observability-essential properties for real-time scheduling in stochastic production environments.


\section{Formal Problem Formulation (Dec-POMDP / Stochastic Game)}
\label{sec:problemformulation}

The dual-resource FJSP with stochastic job arrivals is formulated as a Decentralized Partially Observable Markov Decision Process (Dec-POMDP), which provides a rigorous mathematical framework for multi-agent decision-making under partial observability and uncertainty. This formulation naturally captures the decentralized nature of job-centric scheduling, where each job agent makes resource allocation decisions based on local observations while coordinating implicitly through shared system dynamics and rewards.

\subsection{Mathematical Framework}
A Dec-POMDP is defined by the tuple \( \langle \mathcal{N}, \mathcal{S}, \{\mathcal{A}^i\}, \mathcal{T}, \{\mathcal{O}^i\}, \mathcal{Z}, \mathcal{R}, \gamma \rangle \), where:

\begin{itemize}
    \item \( \mathcal{N} = \{1, 2, \ldots, n(t)\} \) is the set of agents (jobs) at time \( t \), where \( n(t) \) varies dynamically as jobs arrive and complete,
    \item \( \mathcal{S} \in \mathbb{R}^{d_s} \) is the global state space encoding system-wide metrics,
    \item \( \mathcal{A}^i \subseteq \{1, \ldots, M\} \) is the action space for agent \( i \), representing machine selection,
    \item \( \mathcal{T}: \mathcal{S} \times \mathcal{A}^1 \times \cdots \times \mathcal{A}^n \rightarrow \Delta(\mathcal{S}) \) is the state transition function,
    \item \( \mathcal{O}^i \in \mathbb{R}^{d_o} \) is the observation space for agent \( i \),
    \item \( \mathcal{Z}: \mathcal{S} \times \mathcal{A}^1 \times \cdots \times \mathcal{A}^n \rightarrow \Delta(\mathcal{O}^1 \times \cdots \times \mathcal{O}^n) \) is the observation function,
    \item \( \mathcal{R}: \mathcal{S} \times \mathcal{A}^1 \times \cdots \times \mathcal{A}^n \rightarrow \mathbb{R} \) is the shared reward function,
    \item \( \gamma \in [0,1) \) is the discount factor for future rewards.
\end{itemize}

The objective is to learn a joint policy \( \boldsymbol{\pi} = (\pi^1, \ldots, \pi^n) \) that maximizes the expected cumulative discounted reward:
\[
J(\boldsymbol{\pi}) = \mathbb{E}_{\tau \sim \boldsymbol{\pi}} \left[ \sum_{t=0}^{T} \gamma^t r_t \right]
\]
where \( \tau = (s_0, \mathbf{a}_0, r_0, s_1, \mathbf{a}_1, r_1, \ldots) \) denotes a trajectory under policy \( \boldsymbol{\pi} \).

% ===============================
% NEW GROUPED SECTION BEGINS HERE
% ===============================

\subsection{Agent-Level Representation: State, Observation, and Action}

% --- State Space ---
\subsubsection{State Space}

The global state \( s_t \in \mathbb{R}^{10} \) provides a complete snapshot of the system at decision epoch \( t \):
\begin{equation}
s_t = \begin{bmatrix}
n_{\text{jobs}}^{\text{arrived}} & n_{\text{jobs}}^{\text{proc}} & n_{\text{jobs}}^{\text{wait}} & n_{\text{ops}}^{\text{arrived}} & n_{\text{ops}}^{\text{proc}} \\
n_{\text{ops}}^{\text{wait}} & u_M & u_O & W_{\text{total}} & \tau
\end{bmatrix}^\top
\end{equation}
where:
\begin{itemize}
    \item \( n_{\text{jobs}}^{\text{arrived}}, n_{\text{jobs}}^{\text{proc}}, n_{\text{jobs}}^{\text{wait}} \): raw counts of arrived, processing, and waiting jobs,
    \item \( n_{\text{ops}}^{\text{arrived}}, n_{\text{ops}}^{\text{proc}}, n_{\text{ops}}^{\text{wait}} \): raw counts of arrived, processing, and waiting operations,
    \item \( u_M = \frac{\sum_{m=1}^{M} \mathbb{1}[\text{machine } m \text{ busy}]}{M} \in [0,1] \): normalized machine utilization,
    \item \( u_O = \frac{\sum_{k=1}^{K} \mathbb{1}[\text{operator } k \text{ busy}]}{K} \in [0,1] \): normalized operator utilization,
    \item \( W_{\text{total}} = \sum_{j \in J_{\text{active}}} w_j \): cumulative waiting time (raw, unbounded),
    \item \( \tau = \frac{t_{\text{sim}}}{T_{\text{episode}}} \in [0,1] \): normalized episode time fraction.
\end{itemize}

% --- Observation Space ---
\subsubsection{Observation Space}

Each agent \( i \) receives a partial observation \( o_t^i \in \mathbb{R}^8 \) encoding only its own job-specific information:
\begin{equation}
o_t^i = \begin{bmatrix}
\theta_{\text{current}} & n_{\text{ops}}^{\text{total}} & n_{\text{ops}}^{\text{remain}} & w_i & |M_\theta| & |M_{\theta}^{\text{free}}| & n_{\text{jobs}}^{\text{active}} & \delta_{\text{finished}}
\end{bmatrix}^\top
\end{equation}
where:
\begin{itemize}
    \item \( \theta_{\text{current}} \in \{0, 1, \ldots, 8\} \): operation type index of current operation,
    \item \( n_{\text{ops}}^{\text{total}}, n_{\text{ops}}^{\text{remain}} \): total and remaining operations for this job,
    \item \( w_i \geq 0 \): cumulative waiting time for job \( i \) (raw, unbounded),
    \item \( |M_\theta| \): number of machines theoretically capable of operation type \( \theta \) (theoretical capacity),
    \item \( |M_{\theta}^{\text{free}}| \): number of capable machines currently available after feasibility masking,
    \item \( n_{\text{jobs}}^{\text{active}} \): system-wide active job count (shared coordination signal),
    \item \( \delta_{\text{finished}} \in \{0,1\} \): job completion indicator (1 when all operations complete, 0 otherwise).
\end{itemize}

\textbf{Rationale for Active Job Count:} The inclusion of \( n_{\text{jobs}}^{\text{active}} \) in the observation space is critical for enabling congestion-aware learning. Since jobs arrive stochastically (\(\lambda = 0.4\)) and depart upon completion, the agent population \( n(t) \) varies dynamically throughout the episode. When making allocation decisions, each agent must contextualize its local information (\( |M_{\theta}^{\text{free}}| \), \( w_i \)) with the current level of resource competition: a scenario with 3 active jobs and 60\% machine utilization represents fundamentally different decision context than 8 active jobs with 60\% utilization, despite identical resource availability signals. By observing \( n_{\text{jobs}}^{\text{active}} \) alongside resource utilization metrics (machine/operator availability), agents learn to \textit{cluster similar system states} and adapt their policies accordingly. For example, in high-congestion states (\( n_{\text{jobs}}^{\text{active}} \geq 6 \)), agents may prioritize minimizing waiting time over optimizing processing time, whereas in low-congestion states (\( n_{\text{jobs}}^{\text{active}} \leq 2 \)), they can afford to wait for preferred resources. This enables the multi-agent system to develop context-dependent coordination strategies that generalize across varying workload intensities, a key requirement for robust performance in stochastic manufacturing environments.

\textbf{Action Space and Masking Dimensions:} Each agent selects from a discrete action space \( \mathcal{A}^i = \{1, 2, \ldots, M\} \) where \( M = 5 \) machines, resulting in an output dimension of 5 from the agent Q-network. However, action masking \( \mathcal{M}_t^i \in \{0,1\}^5 \) dynamically restricts this space to feasible allocations: a machine is only selectable if it is capable of the required operation, currently idle, \textit{and} at least one qualified operator is available. This dual-resource feasibility check (machine + operator) ensures that selecting a machine implicitly guarantees a valid (machine, operator) pair assignment, avoiding infeasible actions while maintaining a compact action space.

\textbf{Partial Observability:} Each agent observes only its own job state (\( \theta_{\text{current}}, n_{\text{ops}}^{\text{remain}}, w_i \)) and limited system-level information (\( |M_{\theta}^{\text{free}}|, n_{\text{jobs}}^{\text{active}} \)). Agents cannot observe: (1) detailed schedules of other jobs, (2) exact resource allocation timelines, (3) future job arrivals, or (4) the full global state \( s_t \). This enforces realistic information constraints present in distributed manufacturing systems.

% --- Action & Masking ---
\subsubsection{Action Space and Action Masking}

Each agent \( i \)'s action space is a discrete set of machine indices:
\[
\mathcal{A}^i = \{1, 2, \ldots, M\}
\]
where \( M \) is the total number of machines in the system. Although agents formally select machines, \textit{the action masking mechanism ensures that each selectable action implicitly represents a feasible (machine, operator) pair}: agents can only observe and select machines for which both the machine is available \textit{and} a qualified operator exists and is currently free. This implicit pair feasibility is enforced through the three-condition mask validation before agents make decisions.

\textbf{Action Masking:} At each decision epoch, a binary feasibility mask \( \mathcal{M}_t^i \in \{0,1\}^M \) is constructed for agent \( i \):
\begin{equation}
\mathcal{M}_t^i[m] = \mathbb{1}\left[ C_{\theta,m} \land \neg B_m \land \exists k \in \mathcal{K}_m^{\text{free}} : Q_{k,\theta,m} \right]
\end{equation}
where:
\begin{itemize}
    \item \( C_{\theta,m} \): machine \( m \) is capable of operation type \( \theta \) (\textit{capability constraint}),
    \item \( B_m \): machine \( m \) is currently busy processing another operation,
    \item \( \mathcal{K}_m^{\text{free}} \): set of operators currently not assigned to any machine,
    \item \( Q_{k,\theta,m} \): operator \( k \) is qualified for machine \( m \).
\end{itemize}

This three-condition conjunction ensures \textit{dual-resource feasibility}: a machine \( m \) is only selectable (\( \mathcal{M}_t^i[m] = 1 \)) if it is capable, not busy, \textit{and} has at least one qualified operator available. Thus, the masked action space \( \mathcal{A}_t^i = \{ a \in \mathcal{A}^i : \mathcal{M}_t^i[a] = 1 \} \) implicitly represents the set of feasible \textit{(machine, operator)} pairs, where each machine index corresponds to a valid resource combination. If \( \mathcal{A}_t^i = \emptyset \) (either because no capable machines are currently available or because no qualified operator is free for the feasible machines), the agent defers its decision and waits for the next decision epoch.

\textbf{Operator Assignment Coupling:} When agent \( i \) selects machine \( m \) from the masked action space \( \mathcal{A}_t^i \), the environment assigns a qualified operator \( k^* \) from the set of available operators that satisfied the mask condition:
\[
k^* = \arg\min_{k \in \mathcal{K}_m^{\text{free}} : Q_{k,\theta,m}} \, \xi_k, \quad \xi_k \sim \text{Uniform}(0,1)
\]
where the random seed is fixed per episode for reproducibility. Since the action mask already verified \( \exists k \in \mathcal{K}_m^{\text{free}} : Q_{k,\theta,m} \), this assignment is guaranteed to succeed. This design eliminates the combinatorial explosion of explicit \textit{(machine, operator)} action spaces: instead of \( \mathcal{A}^i = \prod_{m=1}^M \prod_{k=1}^K \{(m,k)\} \), the action space remains \( \mathcal{A}^i = \{1,\ldots,M\} \), with masking implicitly encoding pair feasibility. Critically, operators function as an \textit{additional resource constraint} rather than a decision variable.

% ===============================
% TRANSITION + REWARD MERGED HERE
% ===============================

\subsection{Environment Dynamics and Reward Structure}

% --- Transition dynamics ---
\subsubsection{Transition Dynamics and Decision Epochs}

State transitions are governed by a hybrid SimPy discrete-event simulation:
\begin{equation}
s_{t+1} \sim \mathcal{T}(s_t, \mathbf{a}_t), \quad \mathbf{a}_t = (a_t^1, \ldots, a_t^{n(t)})
\end{equation}

\textbf{Decision Epoch Mechanism:} Decision epochs occur when jobs reach decision points in their lifecycle:
\begin{enumerate}
    \item Each active job runs a SimPy coroutine (\texttt{\_job\_process}) that yields control when it requires a resource allocation decision.
    \item The job pushes a decision request to a queue (\texttt{pending\_decisions}) and suspends via \texttt{yield resume\_evt}.
    \item The training loop polls the simulation in discrete time steps (\( \Delta t = 1.0 \)) using \texttt{wait\_for\_decisions()} until \texttt{pending\_decisions} is non-empty.
    \item All pending jobs form a decision batch \( \{j_1, \ldots, j_{|\mathcal{B}_t|}\} \), and agents simultaneously select machines via the policy \( \pi_\theta \).
    \item Selected machines are allocated with environment-assigned operators, operations begin execution, and the simulation resumes until the next decision epoch.
\end{enumerate}

\textbf{Processing Time Heterogeneity and Job Completion Variability:}
Operation processing times \( d_{\theta,m} \) exhibit \textit{processing time variability} across different machines due to varying machine capabilities. For a given operation type \( \theta \), the processing time is \textit{deterministically} determined by the selected machine \( m \), reflecting differences in machine efficiency and technological specifications. These durations are derived from empirically-calibrated processing time means 
\texttt{processing\_time\_means}[\(\theta\)][\(m\)], which capture the expected execution time for each operation-machine combination.

Job completion times are highly \textit{variable} and depend on multiple dynamic factors:
\begin{itemize}
    \item \textbf{Scheduling Decisions:} The sequence and timing of machine-operator allocations directly affect when operations begin and complete.
    \item \textbf{Waiting Times:} Resource contention causes jobs to wait when all feasible machine-operator pairs are busy, accumulating delay that propagates through the job's operation sequence. Since agents cannot directly control operator assignment, operator availability introduces stochastic wait times: a machine may be free but unavailable if no qualified operators are idle.
    \item \textbf{System Load Dynamics:} Stochastic job arrivals following a compound Poisson process \( N(t) \sim \text{Poisson}(\lambda t) \), where \( \lambda \) is the arrival rate parameter, create fluctuating workload conditions. Periods of high arrival density lead to congestion where queue lengths grow and jobs experience extended wait times, while low-arrival periods allow the system to clear backlogs.
    \item \textbf{Dual-Resource Constraint:} Operator availability at each decision epoch depends on current operator utilization. While agents learn to select machines, the environment controls operator assignment, creating an implicit coupling where the same machine selection may yield different wait times depending on operator availability patterns. Agents learn temporal coordination through global operator utilization signals (\( u_O \)) and wait time feedback, adapting machine selection timing to avoid contested resources.
\end{itemize}

Between decision epochs, the simulation advances as operations execute for their prescribed durations \( d_{\theta,m} \), and the global state evolves continuously. New decision epochs are triggered when:
\begin{enumerate}
    \item An operation completes, releasing machine-operator resources and allowing the job to request allocation for its next operation.
    \item A new job arrives stochastically, instantiating a new agent \( i_{\text{new}} \) that immediately enters the decision queue.
\end{enumerate}

\subsubsection{Reward Function: Multi-Objective Optimization}

All agents share a global reward signal designed to balance competing scheduling objectives:
\begin{equation}
r_t = w_1 C_{\text{norm}} - w_2 W_{\text{norm}} - w_3 P_{\text{norm}} + w_4 \Delta_{\text{throughput}} + w_5 B_{\text{load}}
\end{equation}
where:
\begin{itemize}
    \item \( C_{\text{norm}} = \frac{|\{j \in J : j \text{ completed}\}|}{|J|} \in [0,1] \): normalized completion rate,
    \item \( W_{\text{norm}} = \frac{\bar{w}}{W_{\max}} \in [0,1] \): normalized average waiting time, \( \bar{w} = \frac{1}{|J_{\text{active}}|} \sum_{j \in J_{\text{active}}} w_j \),
    \item \( P_{\text{norm}} = \frac{|J_{\text{active}}|}{|J|} \in [0,1] \): normalized work-in-progress ratio,
    \item \( \Delta_{\text{throughput}} = \frac{|\{j : j \text{ completed in last step}\}|}{|J|} \): instantaneous throughput change,
    \item \( B_{\text{load}} = \lambda_M \cdot H(M) + \lambda_O \cdot H(O) \in [0,1] \): entropy-based load balance score,
\end{itemize}

\textbf{Reward Weight Configuration:}
The weights \( w_1, w_2, w_3, w_4, w_5 \) and parameters \( \lambda_M, \lambda_O, W_{\max} \) are hyperparameters that control the relative importance of each objective component.

% ===============================
% POLICY (UNCHANGED)
% ===============================

\subsection{Policy Representation and Learning Objective}

Each agent \( i \) follows a parameterized policy \( \pi_\theta^i(a^i \mid h_t^i) \), where \( h_t^i \) is the agent's recurrent hidden state encoding its observation history. Under the MASA-QMIX framework (detailed in Section~\ref{sec:masaqmix}), agents share RNN parameters but maintain separate hidden states, enabling:
\begin{itemize}
    \item \textit{Parameter Sharing:} Single RNN \( f_\theta: \mathcal{O}^i \times \mathcal{H} \rightarrow \mathbb{R}^{|\mathcal{A}^i|} \times \mathcal{H} \) processes all agents,
    \item \textit{Dynamic Population Handling:} New agents use the same \( f_\theta \) with zero-initialized \( h_0^i \),
    \item \textit{Action Selection:} \( \epsilon \)-greedy exploration:
    \[
    a_t^i = 
    \begin{cases}
    \text{uniform}(\mathcal{A}_t^i) & \text{with probability } \epsilon(t) \\
    \arg\max_{a \in \mathcal{A}_t^i} Q_\theta(h_t^i, a) & \text{otherwise}
    \end{cases}
    \]
    where \( \epsilon(t) \) decays linearly over the training horizon.
\end{itemize}

The learning objective is to minimize:
\[
\mathcal{L}(\theta) = \mathbb{E}_{\tau \sim \mathcal{D}} \left[ \left( Q_{\text{tot}}(\mathbf{s}_t, \mathbf{a}_t; \theta) - y_t \right)^2 \right]
\]
where \( y_t = r_t + \gamma \max_{\mathbf{a}'} Q_{\text{tot}}(\mathbf{s}_{t+1}, \mathbf{a}'; \theta^-) \).


\section{System Architecture}
\label{sec:systemarchitecture}
This section describes the implementation architecture of the MASA-QMIX framework, detailing how the Dec-POMDP formulation (Section~\ref{sec:formulation}) is realized through three integrated components: (1) a SimPy-based discrete-event simulation environment, (2) a recurrent neural network agent architecture with parameter sharing, and (3) the QMIX value decomposition algorithm adapted for dynamic job-centric scheduling.

\subsection{Environment Design: SimPy Discrete-Event Simulation}
\label{sec:env_design}

The scheduling environment is implemented as a discrete-event simulation (DES) using SimPy, a process-based discrete-event simulation framework. This design naturally models asynchronous job arrivals, concurrent operation processing, and resource contention without requiring discretization of continuous time.

\textbf{Core Components:}

\begin{itemize}
\item \textbf{SimPy Environment (\texttt{simpy.Environment}):} Manages simulation time advancement and event scheduling. Unlike traditional RL environments with fixed time steps, SimPy advances time only when events occur (job arrivals, operation completions, decision requests), enabling efficient simulation of sparse scheduling decisions.

\item \textbf{Job Agents (\texttt{JobAgent}):} Each job is represented as a SimPy process (coroutine) that executes a lifecycle: arrival → decision request → resource allocation → operation execution → completion. Jobs run concurrently, yielding control when waiting for decisions or resources.

\item \textbf{Resource Management:} 
\begin{itemize}
    \item \textit{Machine Resources:} Implemented as \texttt{simpy.Resource} objects with capacity 1, enforcing the machine disjunction constraint (one operation per machine at a time).
    \item \textit{Operator Pool:} Managed through a custom \texttt{OperatorPool} class that tracks operator qualifications, availability, and assignment history. Operators are assigned via:
     \texttt{find\_free\_operator\_for\_machine\_seeded\_random()}, which filters qualified free operators and selects one deterministically (seeded random tie-breaking for reproducibility).
\end{itemize}

\item \textbf{Decision Epoch Mechanism:} The bridge between SimPy's event-driven execution and MARL's episodic structure:
\begin{enumerate}
    \item Jobs requiring decisions push requests to a \texttt{pending\_decisions} queue and suspend via \texttt{yield resume\_evt}.
    \item The training loop polls \texttt{wait\_for\_decisions()} in discrete intervals (\(\Delta t = 1.0\)), which advances the simulation until \texttt{pending\_decisions} contains at least one job or the episode terminates.
    \item All pending jobs form a decision batch, and the policy selects actions for the entire batch simultaneously (batch-parallel decision-making).
    \item Actions are validated against the current action mask, allocated to resources, and operations begin execution.
    \item The simulation resumes until the next decision epoch.
\end{enumerate}

\item \textbf{Stochastic Job Arrivals:} Jobs arrive following a Poisson process with rate \(\lambda\). The \texttt{TaskGenerator} implements exponential inter-arrival times: after each job arrival, the time until the next arrival is sampled from an exponential distribution with mean \(1/\lambda\), creating realistic stochastic arrival patterns. The generator runs as a concurrent SimPy process, injecting new jobs into the system at sampled times throughout the episode.
\end{itemize}

\textbf{Observation and State Construction:}
The environment implements canonical builders (\texttt{build\_agent\_obs()}, \texttt{build\_state\_vector()}) in \texttt{utils/env\_obs.py}:
\begin{itemize}
    \item \textbf{Agent Observations:} Extracted directly from \texttt{JobAgent} attributes (current operation type, remaining operations, wait time) combined with real-time system queries (free machine count, active job count).
    \item \textbf{Global State:} Aggregated from environment counters (jobs arrived/processing/waiting, operations arrived/processing/waiting), resource utilization (machine and operator), cumulative wait time, and episode progress fraction.
\end{itemize}

\textbf{Reward Computation:}

The reward function (Eq.~4.5) is computed at each decision epoch by querying current system metrics:
\begin{itemize}
    \item \textit{Completion rate:} Ratio of finished jobs to total arrived jobs.
    \item \textit{Average wait time:} Mean cumulative wait time across active jobs, normalized by \texttt{max\_wait\_time}.
    \item \textit{Work-in-progress:} Ratio of active jobs to total jobs.
    \item \textit{Throughput delta:} Number of jobs completed since the last decision epoch.
    \item \textit{Load balance:} Entropy of machine and operator selection distributions over a sliding window of recent decisions.
\end{itemize}

Rewards are normalized using Welford's online algorithm (running mean and variance) to maintain scale-invariance across different job arrival rates and processing time distributions.

\subsection{Agent Architecture: Recurrent Neural Networks with Parameter Sharing}
\label{sec:agent_arch}

Agent policies are implemented as recurrent neural networks (RNNs) to handle partial observability and variable-length observation histories. The architecture employs parameter sharing across all agents to enable generalization to dynamic agent populations.

\textbf{Network Architecture (\texttt{RNNAgent}):}

\begin{enumerate}
\item \textbf{Input Layer:} Each agent's input at time \(t\) is a concatenated vector:
\[
\mathbf{x}_t^i = [\mathbf{o}_t^i; \mathbf{u}_{t-1}^i; \mathbf{e}^i]
\]
where:
\begin{itemize}
    \item \(\mathbf{o}_t^i \in \mathbb{R}^8\): agent observation (Section~\ref{sec:formulation}).
    \item \(\mathbf{u}_{t-1}^i \in \{0,1\}^M\): one-hot encoding of the last action (if \texttt{last\_action=True}).
    \item \(\mathbf{e}^i \in \{0,1\}^{n(t)}\): one-hot agent ID (if \texttt{reuse\_network=True}), enabling the shared network to distinguish between agents when necessary.
\end{itemize}
Thus, the input dimension is \(\text{input\_dim} = 8 + M + n(t)\) (in practice, \(8 + 5 + 10 = 23\) with default settings).

\item \textbf{Encoder Layer:} A fully connected layer with ReLU activation projects the input to a hidden representation:
\[
\mathbf{h}_{\text{enc}}^i = \text{ReLU}(\mathbf{W}_{\text{fc1}} \mathbf{x}_t^i + \mathbf{b}_{\text{fc1}})
\]
where \(\mathbf{h}_{\text{enc}}^i \in \mathbb{R}^{64}\) (default \texttt{rnn\_hidden\_dim=64}).

\item \textbf{Recurrent Layer:} A single-layer Gated Recurrent Unit (GRU) maintains temporal state:
\[
\mathbf{h}_t^i = \text{GRU}(\mathbf{h}_{\text{enc}}^i, \mathbf{h}_{t-1}^i)
\]
The GRU accumulates observation history and enables the agent to reason about temporal patterns (e.g., recurring resource contention, job arrival bursts).

\item \textbf{Action-Value Head:} A final fully connected layer outputs Q-values for all actions:
\[
\mathbf{Q}_t^i = \text{ReLU}(\mathbf{W}_{\text{out}} \mathbf{h}_t^i + \mathbf{b}_{\text{out}}) \in \mathbb{R}^M
\]
where \(Q_t^i[m]\) represents the estimated Q-value for selecting machine \(m\).
\end{enumerate}

\textbf{Parameter Sharing:} All agents share the same network parameters \(\theta\) but maintain separate hidden states \(\mathbf{h}_t^i\). When a new job arrives, its agent is initialized with \(\mathbf{h}_0^i = \mathbf{0}\) and immediately begins using the current policy without retraining. This enables zero-shot generalization to new agents and ensures scalability to arbitrary agent populations.

\textbf{Action Selection:} During training, agents follow an \(\epsilon\)-greedy policy:
\[
a_t^i = 
\begin{cases}
\text{uniform}(\mathcal{A}_t^i) & \text{with probability } \epsilon(t) \\
\arg\max_{a \in \mathcal{A}_t^i} Q_t^i[a] & \text{otherwise}
\end{cases}
\]
where \(\mathcal{A}_t^i\) is the masked action space (only feasible machines). During evaluation, \(\epsilon = 0\) (greedy exploitation).

The exploration rate \(\epsilon(t)\) decays linearly over cumulative SimPy simulation time:
\[
\epsilon(t) = \max\left(\epsilon_{\text{end}}, \epsilon_{\text{start}} - \frac{t_{\text{sim}}}{T_{\text{anneal}}} (\epsilon_{\text{start}} - \epsilon_{\text{end}})\right)
\]
where \(T_{\text{anneal}} = 0.4 \times n_{\text{epochs}} \times n_{\text{episodes}} \times \text{episode\_limit}\) (default: anneal over the first 40\% of total training time). This SimPy-time-based decay adapts automatically to changes in episode length or training duration.

\section{MASA-QMIX Algorithm}
\label{sec:masaqmix}

This section presents the QMIX value decomposition algorithm and its adaptation to the dynamic job-centric scheduling problem. We first establish the theoretical foundation (value factorization, monotonicity, Individual-Global-Max principle), then describe the implementation-specific adaptations for variable agent populations and dynamic scheduling.

\subsection{QMIX Framework and Adaptation to Job Scheduling}
\label{sec:qmix_framework}

\subsubsection{Theoretical Foundation: Value Decomposition and Monotonicity}

The QMIX algorithm \citep{rashid2018qmix} addresses the credit assignment problem in cooperative multi-agent settings by decomposing the joint action-value function \(Q_{\text{tot}}(\mathbf{s}, \mathbf{a})\) into per-agent utilities \(Q_i(o^i, a^i)\) through a learnable mixing network:

\begin{equation}
Q_{\text{tot}}(\mathbf{s}, \mathbf{a}) = f_{\text{mix}}\left( Q_1(o^1, a^1), \ldots, Q_n(o^n, a^n); \mathbf{s} \right)
\end{equation}

\textbf{Monotonicity Constraint:} The key innovation is enforcing a monotonicity constraint on the mixing function:
\begin{equation}
\frac{\partial Q_{\text{tot}}}{\partial Q_i} \geq 0, \quad \forall i \in \{1, \ldots, n\}
\end{equation}

This constraint ensures the \textit{Individual-Global-Max (IGM)} property: 
\[
\arg\max_{\mathbf{a}} Q_{\text{tot}}(\mathbf{s}, \mathbf{a}) = \left( \arg\max_{a^1} Q_1(o^1, a^1), \ldots, \arg\max_{a^n} Q_n(o^n, a^n) \right)
\]

\textbf{Practical Implication:} IGM enables fully decentralized execution during deployment. Since the joint optimal action is the concatenation of individual greedy actions, each agent can independently select \(a^i = \arg\max_{a} Q_i(o^i, a)\) without requiring communication, access to the global state \(\mathbf{s}\), or knowledge of other agents' actions. This is critical for real-time scheduling where centralized coordination would introduce communication overhead and computational bottlenecks.

\textbf{Centralized Training, Decentralized Execution (CTDE):} During training, the mixer network accesses the global state \(\mathbf{s}\) to compute \(Q_{\text{tot}}\) and learn optimal mixing weights. At deployment, only individual agent networks \(Q_i(o^i, a^i)\) are required, enabling decentralized execution without state sharing.

\subsubsection{Implementation Adaptations for Dynamic Job Scheduling}

The MASA-QMIX implementation adapts the standard QMIX algorithm to handle the unique challenges of dynamic job-centric scheduling:

\textbf{Key Adaptations for Dynamic Job Scheduling:}

\begin{itemize}
\item \textbf{Variable Agent Population Handling:} Unlike standard QMIX scenarios with constant team size, the MASA scheduler must handle \(n(t) \in [1, n_{\max}]\) active jobs at each decision epoch, where \(n_{\max}\) represents the maximum agent capacity configured via \texttt{args.n\_agents}. The system accommodates variable populations through:
\begin{itemize}
    \item \textit{Parameter sharing:} A single agent network \(Q_i(o^i, a^i; \theta)\) processes all agents using shared parameters, enabling zero-shot generalization to newly arrived jobs without retraining or architectural modifications.
    
    \item \textit{Fixed-capacity mixing with zero-padding:} The mixer network architecture is dimensioned for \(n_{\max}\) agents. When \(n(t) < n_{\max}\), the agent Q-value tensor is zero-padded: \(\mathbf{Q}_{\text{padded}} = [Q_1, \ldots, Q_{n(t)}, 0, \ldots, 0] \in \mathbb{R}^{n_{\max}}\). Due to the monotonicity constraint (\(\mathbf{W}_1(\mathbf{s}) = |\text{hyper\_w\_1}(\mathbf{s})| \geq 0\)), zero-padded entries contribute exactly zero to the mixed value:
    \[
    Q_{\text{tot}} = \sum_{i=1}^{n(t)} w_i(\mathbf{s}) Q_i + \sum_{i=n(t)+1}^{n_{\max}} w_i(\mathbf{s}) \cdot 0 = \sum_{i=1}^{n(t)} w_i(\mathbf{s}) Q_i
    \]
    This naturally masks inactive agent slots without explicit masking operations, preserving the IGM property for the active agent subset.
    
    \item \textit{Why QMIX requires fixed capacity:} The hypernetwork architecture inherently requires fixed output dimensions: \(\text{hyper\_w\_1}: \mathbb{R}^{d_s} \rightarrow \mathbb{R}^{d_{\text{mix}} \times n_{\max}}\). This is not a limitation of the implementation but a consequence of QMIX's design—the monotonicity constraint \(\frac{\partial Q_{\text{tot}}}{\partial Q_i} \geq 0\) is enforced through non-negative weight matrices of predetermined size. Truly variable-size mixing would require abandoning either (1) the hypernetwork approach (e.g., using attention mechanisms), or (2) the strict monotonicity guarantee (sacrificing IGM). Since neural network layers cannot dynamically change output dimensions during execution without recompilation, the \(n_{\max}\) constraint is fundamental to all hypernetwork-based value decomposition methods (QMIX, VDN with hypernetworks, Weighted QMIX).
    
    \item \textit{Practical equivalence to dynamic architectures:} When \(n_{\max}\) is configured sufficiently above the expected peak agent count (e.g., \(n_{\max} = 10\) for \(\lambda = 0.5\) jobs/time-unit with \(\text{episode\_limit} = 50\), yielding expected peak \(\approx 5\text{--}6\) agents), the fixed-capacity approach behaves equivalently to fully dynamic architectures—no job is rejected due to capacity constraints, and zero-padding introduces no computational overhead beyond unused matrix elements. This design choice reflects real-world manufacturing constraints where production systems operate within planned capacity limits: facilities periodically close and reopen job intake based on workload, naturally imposing bounded concurrency similar to \(n_{\max}\). Alternative architectures such as attention-based mixers \citep{wang2020qplex} or transformer-based aggregation could eliminate capacity constraints but sacrifice monotonicity guarantees (and thus IGM) while introducing additional architectural complexity and potential training instability—tradeoffs deemed unnecessary given the controlled arrival rates and capacity planning inherent to manufacturing scheduling contexts.
\end{itemize}

\item \textbf{Zero-Shot Generalization:} When a new job arrives mid-episode:
\begin{itemize}
    \item The job instantiates as a new agent with ID \(i_{\text{new}}\).
    \item The agent initializes its RNN hidden state: \(\mathbf{h}_0^{i_{\text{new}}} = \mathbf{0} \in \mathbb{R}^{64}\).
    \item The agent immediately uses the current shared policy \(\pi_\theta\) without retraining.
    \item Since all agents share network parameters \(\theta\), the new agent benefits from the learned scheduling strategy accumulated from previous episodes.
\end{itemize}

\item \textbf{Episode Structure and Termination:} Episodes terminate when:
\begin{itemize}
    \item The SimPy simulation time reaches \texttt{episode\_limit} (default: 50 decision epochs), \textit{or}
    \item A manual termination signal is received (e.g., training interruption).
\end{itemize}
\textit{Important:} Episodes do \textbf{not} wait for all jobs to complete. Due to stochastic arrivals, jobs may still be active when the episode limit is reached. The system operates continuously, with episodes serving as training segments rather than natural completion boundaries. The number of active agents \(n(t)\) fluctuates throughout each episode based on arrival rate \(\lambda\) and processing dynamics.

\item \textbf{Partial Observability Handling:} Agents observe only local job-specific information \(o^i \in \mathbb{R}^8\) (Section~\ref{sec:formulation}), not the full global state \(\mathbf{s} \in \mathbb{R}^{10}\). The RNN architecture (Section~\ref{sec:agent_arch}) maintains temporal context through hidden states \(\mathbf{h}_t^i\), enabling agents to infer system trends (e.g., increasing resource contention, operator bottlenecks) from observation sequences.
\end{itemize}

\subsection{Mixer Network Architecture}
\label{sec:mixer_arch}

The mixer network (\texttt{MARL/network/qmix\_net.py}, class \texttt{QMixNet}) implements a hypernetwork architecture that generates state-dependent mixing weights. The design realizes the monotonicity constraint while maintaining expressiveness through learned non-linear transformations.

\subsubsection{Theoretical Design Rationale}

\textbf{Why Hypernetworks?} Directly parametrizing \(Q_{\text{tot}}(\mathbf{s}, \mathbf{a})\) as a fixed neural network would limit the mixer's ability to adapt to different system states. Instead, QMIX uses \textit{hypernetworks}—small neural networks that generate the parameters (weights and biases) of the mixing network \textit{conditioned on the global state} \(\mathbf{s}\). This enables the mixer to adjust coordination strategies dynamically: different system states (e.g., high vs. low resource utilization) produce different mixing weights, allowing context-aware credit assignment.

\textbf{Enforcing Monotonicity:} To satisfy \(\frac{\partial Q_{\text{tot}}}{\partial Q_i} \geq 0\), the hypernetworks generate only \textit{non-negative} weights. This is enforced by applying an absolute value activation: \(\mathbf{W}_1(\mathbf{s}) = |\text{hyper\_w\_1}(\mathbf{s})|\). The absolute value ensures \(\mathbf{W}_1 \geq 0\) element-wise, which guarantees monotonicity in the mixing operation.

\subsubsection{Network Specifications and Implementation}
\begin{itemize}
\item \textbf{Input Dimensions:} Agent Q-values \(\mathbf{Q} \in \mathbb{R}^{n(t)}\) and global state \(\mathbf{s} \in \mathbb{R}^{10}\).
\item \textbf{Mixing Dimension:} \(d_{\text{mix}} = 32\) (configurable via \texttt{args.qmix\_hidden\_dim}).
\item \textbf{Hypernetwork Type:} Single-layer hypernetworks (default \texttt{args.two\_hyper\_layers=False}). Each hypernetwork is a simple linear transformation: \(\text{HyperNet}(\mathbf{s}) = \mathbf{W}_{\text{out}} \cdot \mathbf{s} + \mathbf{b}_{\text{out}}\).
\item \textbf{Activation:} ELU activation in the first mixing layer, no activation in the second layer (scalar output).
\end{itemize}

\textbf{Forward Pass Implementation:}

The mixing operation proceeds in two layers:

\textbf{Layer 1 (Agent Aggregation):} Combines per-agent Q-values into a hidden representation:
\[
\mathbf{h}_{\text{mix}} = \text{ELU}\left( \mathbf{W}_1(\mathbf{s}) \cdot \mathbf{Q} + \mathbf{b}_1(\mathbf{s}) \right) \in \mathbb{R}^{d_{\text{mix}}}
\]
where:
\begin{itemize}
    \item \(\mathbf{W}_1(\mathbf{s}) = |\text{hyper\_w\_1}(\mathbf{s})| \in \mathbb{R}^{d_{\text{mix}} \times n(t)}\): State-dependent weight matrix. The absolute value (\texttt{torch.abs}) enforces non-negative weights, satisfying the monotonicity constraint.
    \item \(\mathbf{b}_1(\mathbf{s}) = \text{hyper\_b\_1}(\mathbf{s}) \in \mathbb{R}^{d_{\text{mix}}}\): State-dependent bias vector.
\end{itemize}

\textbf{Layer 2 (Scalar Projection):} Projects the hidden representation to a single total Q-value:
\[
Q_{\text{tot}}(\mathbf{s}, \mathbf{a}) = \mathbf{W}_2(\mathbf{s})^\top \cdot \mathbf{h}_{\text{mix}} + b_2(\mathbf{s})
\]
where:
\begin{itemize}
    \item \(\mathbf{W}_2(\mathbf{s}) = |\text{hyper\_w\_2}(\mathbf{s})| \in \mathbb{R}^{d_{\text{mix}}}\): Non-negative state-dependent weight vector.
    \item \(b_2(\mathbf{s}) = \text{hyper\_b\_2}(\mathbf{s}) \in \mathbb{R}\): State-dependent scalar bias.
\end{itemize}

\textbf{Implementation Details:}
\begin{itemize}
\item The implementation handles variable batch sizes \((B, T, n(t))\) by flattening batch and time dimensions: \((B, T, n) \rightarrow (B \cdot T, n)\), performing mixing, then reshaping back to \((B, T, 1)\).
\item Batch matrix multiplication (\texttt{torch.bmm}) processes all timesteps simultaneously for efficiency.
\item The \texttt{two\_hyper\_layers} option (when enabled) adds a ReLU-activated hidden layer to each hypernetwork, increasing expressiveness at the cost of additional parameters. In practice, this is disabled for MASA scheduling as single-layer hypernetworks provide sufficient capacity.
\end{itemize}

\subsection{Cooperation Mechanism}
\label{sec:cooperation}

The mixer network induces implicit coordination among job agents without explicit communication. This section explains how the QMIX architecture realizes cooperative behavior through three interconnected mechanisms.

\subsubsection{Theoretical Basis: Credit Assignment via Gradient Flow}

During training, the TD-error gradient \(\nabla_\theta \mathcal{L}(\theta)\) flows from the scalar \(Q_{\text{tot}}\) back through the mixer to each agent's Q-network:
\[
\frac{\partial \mathcal{L}}{\partial \theta_i} = \frac{\partial \mathcal{L}}{\partial Q_{\text{tot}}} \cdot \frac{\partial Q_{\text{tot}}}{\partial Q_i} \cdot \frac{\partial Q_i}{\partial \theta_i}
\]

The term \(\frac{\partial Q_{\text{tot}}}{\partial Q_i}\) acts as a \textit{credit assignment weight}: it determines how much the change in agent \(i\)'s Q-value contributes to the overall joint value. Agents whose actions significantly affect system-wide performance receive stronger learning signals, accelerating convergence toward cooperative policies.

\textbf{Factorized vs. Centralized Value Functions:} Without factorization, a single centralized Q-network \(Q_{\text{cent}}(\mathbf{s}, \mathbf{a})\) would struggle with exponential action spaces (\(|\mathcal{A}|^n\)) and provide no mechanism for decentralized execution. QMIX factorizes the joint value while preserving representational capacity through the non-linear mixer, enabling both efficient learning and decentralized deployment.

\subsubsection{Implementation Mechanisms}

\textbf{1. State-Dependent Mixing Weights:}

The weights \(\mathbf{W}_1(\mathbf{s})\) and \(\mathbf{W}_2(\mathbf{s})\) vary based on the current system state \(\mathbf{s}\), enabling context-aware coordination. For example:
\begin{itemize}
    \item \textit{Low utilization states} (\(u_M, u_O \ll 1\)): When resources are abundant, the mixer may assign relatively uniform weights to agent Q-values, allowing agents to optimize individual job completion without strong coordination.
    \item \textit{High utilization states} (\(u_M, u_O \approx 1\)): When resources are scarce, the mixer adjusts weights to emphasize coordination, potentially downweighting actions that would exacerbate contention.
    \item \textit{High wait time states} (\(W_{\text{total}}\) large): The mixer learns to prioritize agents (jobs) with longer wait times by upweighting their Q-values, implementing an implicit fairness mechanism.
\end{itemize}

These weight adjustments are \textit{learned}, not hand-coded—the mixer discovers effective coordination strategies through TD-error minimization during training (Section~\ref{sec:training_strategy}).

\textbf{2. Credit Assignment:}

During backpropagation, the gradient \(\frac{\partial Q_{\text{tot}}}{\partial Q_i}\) flows back to each agent's network, allocating credit proportionally to the agent's contribution to the joint return. Agents whose machine selections significantly impact system-wide metrics (e.g., reducing total wait time through load balancing) receive stronger learning signals, accelerating convergence toward cooperative policies.

\textbf{3. Implicit Load Balancing:}

The entropy-based load balance term in the reward function (Eq.~4.5, Section~\ref{sec:formulation}) incentivizes diverse machine selection across the agent population. The mixer learns to coordinate agents to distribute workload evenly, even though individual agents only observe local information and cannot directly perceive other agents' actions. This coordination emerges from the mixer observing trajectories where concentrated machine selection (low entropy) correlates with low \(Q_{\text{tot}}\) due to resource contention.


\section{Training Strategy and Implementation}
\label{sec:training_strategy}

This section details the training procedure for the MASA-QMIX algorithm, covering the episode collection loop, experience replay mechanism, target network synchronization, loss computation, exploration strategy, and hyperparameter configuration.

\subsection{Training Loop and Episode Structure}
\label{sec:training_loop}

The training process follows a cyclical structure alternating between data collection (rollout) and policy optimization (learning):

\begin{enumerate}
\item \textbf{Initialization:} Networks (\texttt{eval\_rnn}, \texttt{target\_rnn}, \texttt{eval\_mixer}, \texttt{target\_mixer}) are instantiated and moved to the configured device (CPU or CUDA). The replay buffer is initialized with capacity for storing complete episodes. RMSprop optimizer is configured with learning rate \(\alpha = 2 \times 10^{-4}\).

\item \textbf{Episode Collection (Rollout Phase):} The \texttt{RolloutWorker} (\texttt{MARL/common/rollout.py}) interfaces with the SimPy environment to collect trajectories:
\begin{itemize}
    \item At each decision epoch, pending jobs form a decision batch \(\{j_1, \ldots, j_{|\mathcal{B}_t|}\}\).
    \item The worker constructs batched observations, applies action masks, and queries the policy network for Q-values.
    \item Actions are selected via \(\epsilon\)-greedy (Section~\ref{sec:exploration}), executed in the environment, and transitions are recorded.
    \item When an episode terminates (episode\_limit reached or manual stop), the full trajectory is stored in the replay buffer as a structured episode dictionary containing keys: \texttt{o}, \texttt{o\_next}, \texttt{u}, \texttt{u\_onehot}, \texttt{avail\_u}, \texttt{avail\_u\_next}, \texttt{r}, \texttt{state}, \texttt{state\_next}, \texttt{terminated}, \texttt{filled}.
\end{itemize}

\item \textbf{Learning Phase:} After collecting a configured number of episodes (default: 1 episode per training iteration), the policy performs gradient updates:
\begin{itemize}
    \item Sample a mini-batch of \(B\) episodes from the replay buffer.
    \item Compute TD-error and loss (Section~\ref{sec:loss_function}).
    \item Backpropagate gradients through agent networks and mixer.
    \item Clip gradients to maximum norm \(g_{\max} = 10.0\) to prevent exploding gradients.
    \item Update network parameters via RMSprop.
    \item Increment training step counter.
\end{itemize}

\item \textbf{Target Network Update:} Every \(C = 200\) training steps, target networks are synchronized with evaluation networks via hard copy: \(\theta^- \leftarrow \theta\) (Section~\ref{sec:target_updates}).

\item \textbf{Checkpoint Saving:} Model checkpoints are periodically saved to \texttt{./MARL/model/qmix/masa\_schedule/} containing agent network and mixer parameters, optimizer state, training step count, and epsilon decay progress.
\end{enumerate}

\textbf{Episode Structure and Padding:} Episodes vary in length \(T \in [1, T_{\max}]\) due to stochastic arrivals and termination conditions. During batch sampling, episodes are zero-padded to the maximum length within the batch, and a binary \texttt{filled} mask \(\mathbf{m}_t \in \{0,1\}\) tracks valid timesteps. The loss function applies this mask to exclude padded transitions from gradient computation.

\subsection{Experience Replay Buffer}
\label{sec:replay_buffer}

The replay buffer (\texttt{MARL/common/replay\_buffer.py}) implements episodic storage to break temporal correlations and improve sample efficiency.

\subsubsection{Episode Storage}

Episodes are stored as complete trajectories rather than individual transitions, preserving temporal structure required for RNN hidden state propagation. Each episode is represented as a dictionary of numpy arrays with time \(T\) as the first dimension:

\begin{itemize}
\item \texttt{o}: Observations, shape \((T, n_{\max}, d_o)\)
\item \texttt{u}: Actions, shape \((T, n_{\max})\)
\item \texttt{r}: Rewards, shape \((T, 1)\)
\item \texttt{state}: Global states, shape \((T, d_s)\)
\item \texttt{terminated}: Episode termination flags, shape \((T, 1)\)
\item \texttt{filled}: Valid timestep mask, shape \((T, 1)\)
\item \texttt{avail\_u}: Action masks, shape \((T, n_{\max}, |\mathcal{A}|)\)
\end{itemize}

When \(n(t) < n_{\max}\), agent-dimension arrays are zero-padded and an implicit agent mask is maintained (though unused due to zero-padding's natural masking via monotonicity, as explained in Section~\ref{sec:qmix_framework}).

The buffer maintains a circular queue with maximum capacity (default: 1000 episodes). When capacity is exceeded, the oldest episodes are discarded (FIFO policy).

\subsubsection{Batch Sampling Strategy}

During training, the \texttt{sample()} method constructs mini-batches:

\begin{enumerate}
\item Randomly sample \(B\) episodes from the buffer (default: \(B = 32\)).
\item Determine the maximum episode length \(T_{\max}\) within the batch.
\item Zero-pad shorter episodes to \(T_{\max}\) and set \(\texttt{filled}[t] = 0\) for padded timesteps.
\item Stack episodes along batch dimension to form tensors with shape \((B, T_{\max}, \ldots)\).
\item Return batch dictionary with all required keys for \texttt{QMIX.learn()}.
\end{enumerate}

This uniform random sampling ensures diverse experience distribution and prevents overfitting to recent trajectories. The episodic structure preserves RNN hidden state continuity: each episode begins with \(\mathbf{h}_0^i = \mathbf{0}\), and hidden states evolve sequentially through the episode without cross-contamination between episodes in the batch.

\subsection{Target Network Updates}
\label{sec:target_updates}

To stabilize Q-learning, MASA-QMIX employs separate target networks (\texttt{target\_rnn}, \texttt{target\_mixer}) that are updated infrequently:

\[
\theta^-_{\text{rnn}} \leftarrow \theta_{\text{rnn}}, \quad \theta^-_{\text{mixer}} \leftarrow \theta_{\text{mixer}}, \quad \text{every } C \text{ training steps}
\]

where \(C = 200\) (configurable via \texttt{args.target\_update\_cycle}). This hard synchronization prevents the moving target problem where Q-values and TD-targets shift simultaneously during gradient updates, which can cause training instability and divergence.

\textbf{Implementation:} The \texttt{QMIX} policy class maintains a \texttt{train\_step} counter incremented after each \texttt{learn()} call. When \texttt{train\_step \% C == 0}\), target network parameters are copied via:
\begin{verbatim}
self.target_rnn.load_state_dict(self.eval_rnn.state_dict())
self.target_mixer.load_state_dict(self.eval_mixer.state_dict())
\end{verbatim}

Between updates, target networks remain frozen, providing stable TD-targets:
\[
y_t = r_t + \gamma \max_{\mathbf{a}'} Q_{\text{tot}}(\mathbf{s}_{t+1}, \mathbf{a}'; \theta^-) \quad \text{(targets computed with } \theta^- \text{)}
\]

\subsection{Loss Function and TD Error}
\label{sec:loss_function}

The policy minimizes the mean squared temporal difference error over valid transitions:

\begin{equation}
\mathcal{L}(\theta) = \frac{1}{\sum_{b,t} m_{b,t}} \sum_{b=1}^{B} \sum_{t=1}^{T} m_{b,t} \left( Q_{\text{tot}}(\mathbf{s}_{b,t}, \mathbf{a}_{b,t}; \theta) - y_{b,t} \right)^2
\end{equation}

where:
\begin{itemize}
\item \(m_{b,t} = \texttt{filled}[b,t] \in \{0,1\}\): validity mask excluding padded transitions,
\item \(Q_{\text{tot}}(\mathbf{s}_{b,t}, \mathbf{a}_{b,t}; \theta)\): Mixed Q-value from evaluation networks,
\item \(y_{b,t} = r_{b,t} + \gamma Q_{\text{tot}}(\mathbf{s}_{b,t+1}, \mathbf{a}'_{b,t+1}; \theta^-) \cdot (1 - \texttt{terminated}_{b,t})\): TD-target.
\end{itemize}

\textbf{Action-Value Computation:}

\begin{enumerate}
\item \textbf{Evaluation Q-values:} Compute individual agent Q-values for the \textit{taken actions} \(\mathbf{a}_{b,t}\):
\[
Q_i(h_{b,t}^i, a_{b,t}^i; \theta) \quad \text{for each agent } i
\]
where \(h_{b,t}^i\) is the RNN hidden state at time \(t\). These are gathered via \texttt{torch.gather()} using action indices \(\mathbf{a}_{b,t}\).

\item \textbf{Target Q-values:} Compute Q-values for \textit{all actions} at the next state:
\[
Q_i(h_{b,t+1}^i, a'; \theta^-) \quad \text{for all } a' \in \mathcal{A}
\]
Apply action masking to invalid actions: \(Q_i(h_{b,t+1}^i, a'; \theta^-) \leftarrow -10^{10}\) if \(\texttt{avail\_u\_next}_{b,t+1}[i,a'] = 0\), ensuring only feasible actions are considered during \(\max\) operation.

\item \textbf{Double-Q Target:} Select greedy actions based on masked Q-values:
\[
\mathbf{a}'_{b,t+1} = \arg\max_{a' \in \mathcal{A}_{b,t+1}} Q_{\text{tot}}(\mathbf{s}_{b,t+1}, \mathbf{a}'; \theta^-)
\]

\item \textbf{Mixing:} Aggregate agent Q-values via mixer networks to obtain scalar \(Q_{\text{tot}}\).
\end{enumerate}

\textbf{Gradient Clipping:} After backpropagation, gradients are clipped to maximum L2 norm \(g_{\max} = 10.0\) via \texttt{torch.nn.utils.clip\_grad\_norm\_(self.params, 10.0)} to prevent exploding gradients common in RNN training.

\subsection{Exploration Strategy (Epsilon-Greedy)}
\label{sec:exploration}

Action selection follows an \(\epsilon\)-greedy policy with time-based linear decay:

\begin{equation}
a_t^i = 
\begin{cases}
\text{uniform}(\mathcal{A}_t^i) & \text{with probability } \epsilon(t_{\text{sim}}) \\
\arg\max_{a \in \mathcal{A}_t^i} Q_i(h_t^i, a; \theta) & \text{otherwise}
\end{cases}
\end{equation}

where \(\mathcal{A}_t^i\) is the masked action space containing only feasible actions.

\textbf{Decay Schedule:} The exploration rate decays linearly based on cumulative SimPy simulation time:

\begin{equation}
\epsilon(t_{\text{sim}}) = \max\left( \epsilon_{\text{end}}, \epsilon_{\text{start}} - \frac{t_{\text{sim}}}{T_{\text{anneal}}} (\epsilon_{\text{start}} - \epsilon_{\text{end}}) \right)
\end{equation}

with default parameters:
\begin{itemize}
\item \(\epsilon_{\text{start}} = 1.0\): Initial exploration (pure random),
\item \(\epsilon_{\text{end}} = 0.05\): Final exploration (5\% random),
\item \(T_{\text{anneal}} = 0.4 \times n_{\text{epochs}} \times n_{\text{episodes}} \times \texttt{episode\_limit}\): Anneal over the first 40\% of total training time.
\end{itemize}

\textbf{SimPy-Time-Based Decay:} Unlike traditional RL where epsilon decays per environment step or episode, MASA-QMIX decays per SimPy simulation time. This accounts for variable episode lengths and decision densities—episodes with more arrivals contribute more to exploration progress, naturally balancing exploration across different load conditions.

\subsection{Optimization and Hyperparameters}
\label{sec:hyperparameters}

\textbf{Optimizer:} RMSprop with learning rate \(\alpha = 2 \times 10^{-4}\). RMSprop's adaptive per-parameter learning rates handle the heterogeneous gradient scales between agent networks (processing local observations) and the mixer (aggregating global value).

\textbf{Network Architecture:}
\begin{itemize}
\item Agent RNN hidden dimension: \(d_{\text{rnn}} = 64\)
\item Mixer hidden dimension: \(d_{\text{mix}} = 32\)
\item Activation functions: ReLU (agent network), ELU (mixer layer 1)
\end{itemize}

\textbf{Training Configuration:}
\begin{itemize}
\item Discount factor: \(\gamma = 0.99\)
\item Batch size: \(B = 32\) episodes
\item Replay buffer capacity: 1000 episodes
\item Target update frequency: \(C = 200\) training steps
\item Gradient clipping: \(g_{\max} = 10.0\)
\item Episode limit: \(T_{\max} = 50\) decision epochs
\item Maximum agent capacity: \(n_{\max} = 10\)
\end{itemize}

\textbf{Reward Function Weights (Eq.~4.5):}
\begin{itemize}
\item \(w_1 = 1.0\): Completion rate reward
\item \(w_2 = 0.5\): Waiting time penalty
\item \(w_3 = 0.3\): Work-in-progress penalty
\item \(w_4 = 0.7\): Throughput delta reward
\item \(w_5 = 0.2\): Load balance reward
\item \(\lambda_M = 0.5\), \(\lambda_O = 0.5\): Machine/operator entropy weights
\item \(W_{\max} = 100.0\): Waiting time normalization constant
\end{itemize}

\textbf{Job Arrival Configuration:}
\begin{itemize}
\item Arrival rate: \(\lambda = 0.5\) jobs per time unit (Poisson process)
\item Job generation: Jobs are generated stochastically at runtime by \texttt{TaskGenerator}. Each job consists of a random sequence of operations selected from the available operation types, with sequence lengths varying based on manufacturing process requirements.
\item Operation types: 9 distinct operation types (Op0--Op8, indices 0--8)
\item Machines: 5 machines (configurable via \texttt{args.n\_actions}, determined by environment's \texttt{machine\_list})
\item Processing times: Deterministic per (operation, machine) pair, derived from \texttt{processing\_time\_means} configuration matrix, typically ranging 5--15 time units depending on machine capability and operation complexity
\end{itemize}

\textbf{Checkpoint and Logging:} Model checkpoints are saved every 50 training epochs to \texttt{./MARL/model/qmix/masa\_schedule/} containing network parameters, optimizer state, epsilon progress, and training step count. Training metrics (episode rewards, loss, TD-error, gradient norms) are logged to \texttt{my\_data\_and\_graph/historydata/} for analysis and visualization.
